\documentclass[a4paper, 12pt]
{article}
\usepackage{amsmath, amssymb, esvect, mathtools}

\newcommand{\im}{\text{im}}
\newcommand{\stab}[1]{\text{stab}_G(#1)}
\newcommand{\inv}[1]{#1^{-1}}
\newcommand{\pgroup}{$p-$ group}
\newcommand{\psubgroup}{$p-$subgroup}

\begin{document}
\textbf{Math 321: Groups and Symmetry}

\hfill

Personal revision notes I made for the 2025-2026 LU Math 321 exams
	
	\hfill
	
\textbf{Chapter 0: Basics}
	\begin{itemize}
		\item Common groups:
		\begin{itemize}
		\item $C_n = \{1,w,w^2, \cdots 
		, w^{n-1}\}$ where $w = e^{2\pi i/n}$ together with multiplication
		
		\item \textbf{Group of integers modulo n} is  $\mathbb{Z}_n = \{\hat{0}, \hat{1}, \cdots, \hat{n-1}\}$ together with addition modulo n.
		
		\item  \textbf{Symmetric group of order n} is $S_n$, the group of permutations of $\{1,2, \cdots, n\}$ together with composition
		
		\item\textbf{ Alternating group of order n} is  $A_n$, the subgroup of $S_n$ that consists of all the even permutations
		
		\item \textbf{Dihedral Group }$D_{2n}= \{b^sa^t: b^2= a^n=e,  ab=ba^{n-1}\}$ together with multiplication
		
		In $D_{2n},$ we have $a^kb= ba^{-k}$ for all $k \in \mathbb{Z}$
		
		\item \textbf{Cyclic group} $<g>$
		
		If $C_n = <g>$ where $g^n = e$ and $m |C_n$ then $C_n$ has a subgroup of order m namely $H= <x^d>$ where $d = n/m$
 		
		\item If $G$ is a group and $N$ is a normal subgroup of $G$, the \textbf{quotient group} $G/N$ (read as "$G$ modulo $N$)"), is the set of all cosets of $N$ in $G$, i.e., $G/N = \{gN: g \in G\}$ where the operation is defined as : $\forall g,h \in G, (gN)(hN) = (gh)N$
		
		\end{itemize}
		
		
		\hfill
		
\textbf{Chapter 1: Conjugacy}

		\item Let G be a group. and $g,h \in G.$ We say that $g$ commutes with $h$ if $gh=hg.$
		
		\item \textbf{Centraliser} The centraliser of $g$ in $G$, $C_G(g),$ is the set of all elements in G which commute with $g.$ 
		
		\item How to calculate centralisers
		\begin{itemize}
			\item Make a Cayley table and compare values
		\end{itemize}
		
		
		\item \textbf{Centre} The centre of $Z(G)$ is defined as $Z(G)= \bigcap_{g \in G}C_G(g).$ It is the set of all elements in $G$ which commutes with all the elements in $G.$
		
		$G$ is abelian iff $Z(G)=G$
		
		$Z(G) \triangleleft G$
		
		In a dihedral group, $G= D_{2n},$ $Z(G)= \{e\}$ when $n$ is odd and $Z(G)= \{e, a^{n/2}\}$ when $n$ is even.
		
		\item \textbf{Conjugate} If $g,h \in G,$ h is conjugate to g if $(\exists x \in G)(h=xgx^{-1})$
		
		Conjugacy is an equivalence relation on G. 
		
		\item \textbf{Conjugacy Clases} The equivalence classes are called conjugacy classes. 
		
		$C_g$ is the conjugacy class containing $g,$ i.e., $$C_g = \{h \in G: (\exists x \in G)( h= xgx^{-1})\}$$
		
		$C_g =\{g\} \iff g \in Z(G).$
		
		Properties preserved under conjugation:
		Order and size of centralisers
		
		If g and h are conjugate and $h= xgx^{-1}$ then $C_G(h)= xC_G(g)x^{-1}$
		
		$|C_g| = |G| / |C_G(g)|$
		
		\item Technique to find the conjugacy classes of a group:
		\begin{itemize}
			\item  Determine $Z(G).$ You can simply find elements which commute with every other element. Prove this!
			
			\item 	Start with the elements in $Z(G).$ Each of these are in a class of its own.
			
			\item Find elements which have the same order. They might be in the same class. 
			
			\item Check size of conjugacy class using  	$|C_g| = |G| / |C_G(g)|.$ Bound it by finding a lower bound for $|C_G(g)|$
			
			\item  Check conjugacy conditions between elements and stop when we have enough elements to meet the size of the conjugacy class 
		\end{itemize}
		
			A normal subgroup is a union of conjugacy classes of G.
		
		\item How to find the normal subgroups of a group
		\begin{itemize}
			\item Find the conjugacy classes
			
			\item Consider all the possible unions of these classes
			
			\item Check it is a subgroup using Lagrange's Theorem.
			
			\item If you multiply any two elements in your union, the result must stay inside the union.
			
			Remember: Any subgroup must contain the identity!!  So the union must always contain $C_e$
		\end{itemize}
		
		\item \textbf{Cycle Type} Let $\pi \in S_n.$ The cycle type type($\pi$) of $\pi$ is the list of lengths of cycles in the full disjoint cycle notation for $\pi$ in order of decreasing size.
		
		Cycle type is preserved under conjugation
		
		\item Two permutations are conjugate iff they have the same cycle type.
		
		Remember: when comparing cycle types, make sure both are written in disjoint cycle notation!!
		
		\item For conjugate elements of $S_n,$  $\pi$ and  $\rho,$ how to find $\sigma$ such that $\rho = \sigma \pi \sigma^{-1}$
		
		\begin{itemize}
			\item Write $\rho$ below $\pi$ ensuring the cycles of the same length are matched against each other. Let $\sigma$ be the permutation sending each element of the first row to the element below it.
		\end{itemize}
		
		\item How do you calculate the size of the conjugacy class $C_{\pi}$ given $\pi \in S_n$?
		
		\begin{itemize}
			\item We need to count the number of permutations having the same cycle type as $\pi$
			
			\item Suppose $\pi$ has $b_i$ cycles of length $i.$ Then
			$$|C_{\pi}| = \frac{n!}{1^{b_1} 2^{b_2} \cdots n^{b_n} b_1! b_2! \cdots b_n!}$$
		\end{itemize}
 	\end{itemize}
 	
 	\hfill
 	
 	\textbf{Chapter 2: Direct Products}
 	\begin{itemize}
 		\item \textbf{External Direct Product} If $G_1$ and $G_2$ are groups, the external direct product is the group $G_1 \times G_2$ with the binary operation $(g_1,g_2)(h_1,h_2)= (g_1h_1, g_2h_2)$ $G_1$ and $G_2$ are called \textbf{factors} of $G_1 \times G_2$
 		
 		\item A direct product is abelian iff its factors are abelian
 		
 		\item $o(g_1, g_2) = lcm(o(g_1), o(g_2))$
 		
 		\item $r$ and $s$ are coprime iff $C_r \times C_s \cong C_{rs}$
 		
 		\item Checking if a direct product is isomorphic to some group G:
 		
 		\begin{itemize}
 			\item Compare the orders of the elements in the direct product and the orders of the elements in G
 			
 			\item Check if both G and the direct product is abelian
 		\end{itemize}
 		\item Finding subgroups of a Direct Product
 		
 		\begin{itemize}
 			\item If $H_1 \le G_1$ and $H_2 \le G_2$ then $H_1 \times H_2 \le G_1 \times G_2.$
 			
 			Note:  The direct product can have  subgroups that is not a direct product of subgroups of the factors  
 		\end{itemize}
 	
 	\item Criterion for a group to be a direct product, i.e., what conditions G must hold for there to be an isomorphism $\theta: G \rightarrow G_1 \times G_2$ such that $\theta(G_i) = \tilde{G_i}?$
 	\begin{itemize}
 		\item G must have normal subgroups $G_1$ and $G_2$ such that $G= G_1G_2$ and $G_1 \cap G_2 = \{e\}$ (i.e., every element of $g$ can be written uniquely as $g= g_1g_2$ with $g_i \in G_i$).
 	\end{itemize} 
 	
 	\item \textbf{Internal Direct Product} If G has normal subgroups $G_1$ and $G_2$ satisfying the criterion, then we call G the internal direct product of $G_1$ and $G_2$ and write $G= G_1 \times G_2.$ 

	\item Finite abelian groups classification
	Every finite abelian group is isomorphic to a direct product of cyclic groups in the form
	$$C_{{p_1}^{r_1}} \times C_{{p_2}^{r_2}} \times \cdots \times C_{{p_t}^{r_t}}$$
	where the $p_i$ are primes (not necessarily distinct) and the product is unique up to reordering of the factors
	
	\item How to find all abelian groups of order n (up to isomorphism)
	
	\begin{itemize}
		\item Factorise n into a product of prime powers
		
		\item Find all possible sets of prime powers
		
		\item If $\{p_1^{r_1}, p_2^{r_2}, \cdots p_t^{r_t}\}$ is a set of prime power then $$C_{{p_1}^{r_1}} \times C_{{p_2}^{r_2}} \times \cdots \times C_{{p_t}^{r_t}}$$ is one of the groups of order n.
	\end{itemize}
	
	\item Standard way of writing a finite abelian group G as a direct product of cyclic groups:
	
	\begin{itemize}
		\item For each distinct prime dividing $|G|,$ write the powers which appear in a row, in ascending order of size, keeping the right hand ends of the various rows aligned. Each prime having its own row.
		
		\item Multiply together the numbers occuring withing each column to form numbers $m_1, m_2, \cdots m_l.$
		
		\item Then $G \cong C_{m_1} \times C_{m_2} \times \cdots \times C_{m_l}$ and $m_i | m_{i+1}$ for all $i \in \{1, \cdots, l-1\}.$
	\end{itemize}
	
	\item Classification of subgroups of finite abelian groups
	
	Finite abelian groups have subgroups of all orders allowed by Lagrange's Theorem, i.e., if $G$ is abelian with order $n$ and $m|n$ then $G$ has a subgroup of order $m.$
	
	\item How to find subgroubs of finite abelian groups
	
	\begin{itemize}
		\item Write $G= C_{{p_1}^{r_1}} \times C_{{p_2}^{r_2}} \times \cdots \times C_{{p_t}^{r_t}}$
		
		\item To find a subgroup of order m, since $m|n$ we write 
		$$m= p_1^{t_1}p_2^{t_2} \cdots p_t^{t_t}$$
		\item for each $i,$ write $C_{p_i^{r_i}} = <x>$ and $p_i^{r_i} = dp_i^{t_i}.$ Then $H_i = <x^d>$ is a subgroup of $C_{p_i^{r_i}}$
		
		\item $H_1 \times H_2 \times \cdots \times H_t$ is a subgroup of G.
	\end{itemize}
	
	\item Condition which forces an abelian group to be cyclic If an abelian group has order n and n is a product of distinct primes then G is cyclic.
 	\end{itemize}
	
\hfill

\textbf{Chapter 3: Jordan- Hölder Theorem}
\begin{itemize}
	\item Facts about normal subgroups of a finite group G
	\begin{itemize}
		\item $G$ abelian $\implies$ all subgroups are normal
		\item $\{e\}$ and $G$ are normal
		\item a subgroup with index 2 is normal
		\item a subgroup which is the union of conjugacy classes is normal
		\item $Z(G)$ is normal
		\item Kernel of a homomorphism is normal
	\end{itemize}
	
	\item \textbf{First Isomorphism Theorem}
	
	If $\phi: G \rightarrow  H$ is  a homomorphism then $G/ \ker \phi \cong \im \phi.$
	
	\item \textbf{Second Isomorphism Theorem}
	
	If $N \triangleleft G$ and $H \leq G$ then $(NH)/ N \cong H/ (N \cap H)$
	
		\item \textbf{Simple group} A non-trivial group $G$ is simple iff it has no proper non-trivial normal subgroups. 
	
	\item Classification of simple abelian groups
	
	Suppose $G$ is abelian. Then $G$ is simple iff $G \cong C_p$ for some prime $p.$
	
	\item \textbf{Maximal} A proper normal subgroup $N \triangleleft G$ is called maximal if $N$ is not contained in any normal subgroup of $G$ except $G$ and itself, i.e, If $M$ is normal and $N \triangleleft M$ then $M= N$ or $M=G.$
	
	\item A normal subgroup $N \triangleleft G$ is maximal iff $G/N$ is simple.
	
	\item If $N \triangleleft G$ satisfies $|G|/|N| = p$ for some prime $p$ then $N$ is maximal and $G/N \cong C_p$
	
	\item \textbf{Composition series} A composition series for G is a series of subgroups
	$$G= G_0 \geq G_1 \geq G_2 \geq \cdots \geq G_r = \{e\}$$ such that $G_i \triangleleft G_{i-1}$ and $G_{i-1}/G_i$ is simple for all i.  The simple groups $G_{i-1}/G_i$ are called the \textbf{composition factors} of G.
	
	\item How to find the composition factor of a group
	\begin{itemize}
		\item Start with a non-trivial group G and choose a maximal normal subgroup, call it $G_1.$ (The trivial subgroup is such a group so a maximal one exists.) You can check $G_1$ is maximal by calculating $|G|/|G_1|.$ 
		\item Repeat the procedure and choose a maximal normal subgroup of $G_i$ for $i \geq 1.$ until we obtain the trivial subgroup
	
	\end{itemize}
	
	\item \textbf{Jordan- Hölder Theorem}
	The list of composition factors for any two composition series is the same
	
	We can associate a collection of simple groups (the composition factors) to a given finite group
	
	\item \textbf{Extension of N by Q} If $G, N,$ and $Q$ are groups with $N \triangle G$ and $G/N \cong Q$ then $G$ is called an extension of $N$ by $Q.$ 

	
	\item \textbf{Soluble} A group $G$ is soluble if its composition factors are all abelian
	
	$G$ abelian $\implies$ G soluble.
	
	\item Sufficient conditions for $N \triangleleft A_n$ to be equal to $A_n$ where $n \geq 5,$ i.e., for $N$ to be a simple group.
	
	\begin{itemize}
		\item If $N$ contains a 3-cycle
		\item If $N$ contains a product of two disjoint transpositions
	\end{itemize}
	
	\item How to find a 3-cycle in $N \triangleleft A_n$ given $\pi \in N$
	
	\begin{itemize}
		\item Write $\pi$ as a product of disjoint cycles $\pi = \gamma_1 \gamma_2 \cdots\gamma_k $in non-increasing length
		
		\item Case 1:  $\gamma_1 = (a_1 a_2 \cdots a_m)$ where $m > 3$. Set $\sigma = (a_1 a_2 a_3)$ and calculate $\sigma \pi \sigma^{-1} \pi^{-1}$. You should get $(a_1 a_2 a_4)$
		
		\item Case 2: $\gamma_1 = (a_1 a_2 a_3)$ and $\gamma_2 = (a_4 a_5 a_6).$ Set $\sigma = (a_2 a_3 a_4)$ and calculate $\sigma \pi \sigma^{-1} \pi^{-1}.$  You should get $(a_1 a_2 a_4 a_3 a_5).$ Apply case 1.
		
		\item Case 3: $\gamma_1 = (a_1 a_2 a_3)$and $ \gamma_2 \gamma_3 \cdots\gamma_k $ are transpositions. Calculate $\pi^2.$ You should get $(a_1 a_3 a_2)$
		
		\item Case 4: $\gamma_1 \gamma_2 \cdots\gamma_k$ are all transpositions. Let $\gamma_1 = (a_1 a_2)$ and $\gamma_1 = (a_3 a_4).$ Set $\sigma = (a_2 a_3 a_4)$ Calculate $\sigma \pi \sigma^{-1} \pi^{-1}$ You should get $(a_1a_4)(a_2a_3)$
	\end{itemize}
	
	\item Conclusion: $A_n$ where $n \geq 5.$ is an example of a non-abelian simple group and hence non-soluble group.
	\end{itemize}
	
	\hfill
	
	\textbf{Chapter 4: Group Actions}
	
	\begin{itemize}
		\item  \textbf{Action} We say that a group G acts on a non-empty set X if for all $g \in G$ and for all $x\in X$ there exists a unique  $gx \in X$ such that 
		
		\begin{align*}
			(g_1g_2)x &= g_1(g_2x)\\
			ex&=x 
		\end{align*}
		
		Note this is a function $G \times X \rightarrow X.$
		
		\item Examples of actions
		\begin{itemize}
			\item \textbf{Rotation} Let $G = \mathbb{R}$ (under addition) and $X = \mathbb{R}^2.$ Define action $G \times X \rightarrow X$ by $g(x,y) \mapsto R_g(x,y)$
			
			\item \textbf{Natural Action} Let $G =S_n$ and $X = \{1,2, \cdots,n\}.$ Define the action $G \times X \rightarrow X$ by $\pi i \mapsto \pi (i)$ for all $\pi \in S_n$ and $i \in X.$
		\end{itemize}
		
	\item \textbf{Orbit} If $G$ acts on $X,$ the orbit of $x \in X$ under $G$ is the set $O_x = \{gx: g \in G\}.$ The set of points to which $x$ can be moved by elements of $G.$
	
	\item If $G$ acts on $X,$ the orbits of the action partition $X.$

 	\item \textbf{Stabiliser} If $G$ acts on $X,$ the stabiliser of $x \in X$ are all the elements in $G$ which fixes $x$ i.e.,
 	$$\stab{x}= \{g \in G: gx=x\}$$
 	
 	\item Stabilisers are subgroups of $G$
 	
 	\item \textbf{Orbit-Stabiliser Theorem} If $G$ acts on $X$ and $x \in X$ then $$|O_x| = |G: \stab{x}|$$
 	
 	If G is is finite, by Lagrange's Theorem $|O_x||\stab{x}|= |G|.$
 	
 	\item  Points lying in the same orbit have conjugate stabilisers, i.e., $$gx=y \implies g\stab{x}\inv{g} = \stab{y}$$
	
	\item \textbf{Polyhedron} Solid shape whose faces are polygons. It is \textbf{regular} if all the faces are regular polygons of the same size and shape with the configuration at each vertex being the same.
	
	\item Five types of regular polyhedron
 	 \end{itemize}
 	
 	\hfill
 	
 	\textbf{Chapter 5: $p$-groups and $p$-subgroups}
 	
 	\begin{itemize}
 		\item \textbf{\pgroup} A group whose order is a positive power of $p$ where $p$ is prime
 		
 		\item \textbf{\psubgroup} A subgroup of $G$ whose order is a power of $p$ 
 		
 		\item \textbf{Sylow's First Theorem} If $G$ is a group of order $p^mr$ where $p \not | r,$ then $G$ has a subgroup of order $p^m$
 		
 		\item \textbf{Sylow \psubgroup} \textit{(Tells us a sufficient condition for a group to have a p-subgroup)}
 		
 		 If $G$ is a group of order $p^mr$ where $p \not | r,$ then a subgroup of order $p^m$ is called a Sylow $\psubgroup$ of $G.$
 		
 		\item \textbf{Sylow's Second Theorem} \textit{(Relates the conjugacy of two Sylow p-subgroups)}
 		
 		If $G$ is a finite group, then any two Sylow p-subgroups $P_1, P_2$ of $G$ are conjugate in $G,$ i.e., there exists $g \in G$ such that $P_2 = gP_1g^{-1}.$
 		
 		\item \textbf{Sylow's Third Theorem} \textit{(gives us information on the number of Sylow p-subgroups)}
 		

 		 If $G$ is a group of order $p^mr$ where $p \not| r$ annd $n_p$ is the number of distinct Sylow p-subgroups of $G$, then $n_p|r$ and $n_p \equiv 1\mod p,$ i.e., $p|(n_p-1)$
 	\end{itemize}
 	
\end{document}